\section{The continuous relaxation and its exact inner-factor defect (NS-76)}
\label{sec:ns76-inner-distance}

\paragraph{Scope and prior work.}
Burnol's quantitative Nyman formula identifies the $q=1$ continuous
approximation defect through the off-critical Blaschke product
\cite[Theorem 1.2 and Lemma 2.4]{BurnolNyman2001}. We apply the same
classical Hardy-space mechanism to this project's $q=2$ target, keeping
its full exterior tail. The calculation gives an exact distance for
\emph{continuous} dilations and a lower bound for integer dilations.
No unconditional equality of the two nonzero distances is assumed.
This is an explicit reformulation of the unknown arithmetic inner
factor, not a new zero-free estimate.

Let $\mathcal T$ be the complete tail space from NS-74, and let
$\mathbb H^2$ denote the Hardy space of $\Re s>1/2$, with boundary norm
$(2\pi)^{-1}\int_{\mathbb R}|f(1/2+it)|^2dt$.
Multiplication on the Mellin boundary by $(s-1)/s$ induces a unitary
map $W$ on the full $L^2$ space. It maps $\mathcal T$ onto
$L^2(0,1)$, whose Mellin image is $\mathbb H^2$.
Indeed in logarithmic coordinates it is
\[
 (Wu)(t)=u(t)-\int_{-\infty}^t e^{-(t-r)/2}u(r)\,dr.
\]
For $t<0$ a tail vector is $C e^{t/2}$ and cancels exactly in this
formula. Conversely the adjoint applied to a function supported on
$t\geq0$ has on $t<0$ a constant multiple of $e^{t/2}$.
Thus the map is onto, and the full tail has been accounted for.

The transformed target and first atom have Mellin transforms
\[
 T(s)=\frac{s-1}{s^3},\qquad
 F(s)=-\frac{(s-1)\zeta(s)}{s^3}.
\]
Write $B$ for the off-critical Blaschke product, normalized by $B(1)>0$:
\[
 B(s)=\prod_{\substack{\zeta(\rho)=0\\\Re\rho>1/2}}
 \frac{s-\rho}{s-(1-\overline\rho)}
 \frac{1-\overline\rho}{\rho}
 \left|\frac{\rho}{1-\rho}\right|.
\]
Multiplicities are included. Burnol's factorization proves that the
inner factor of $(s-1)\zeta(s)/s^2$ is exactly $B$: analytic
continuation excludes a finite-boundary singular factor, and the
real-axis asymptotic excludes an exponential factor. The additional
factor $1/s$ is outer, so $F$ has the same inner factor. The product
and its derivatives converge locally in $\Re s>1/2$ away from its
zeros; $B(1)$ and its first two derivatives are real by conjugation
symmetry.

\begin{proposition}[Exact distance for the complete continuous family]
\label{prop:ns76-continuous-distance}
Let $\mathcal V_{\rm cont}$ be the closed span in the original Hilbert
space of $\sigma_u$, for all real $u\geq1$. Its image under $W$ and
Mellin transformation is $B\mathbb H^2$.
Put $b=B(1)$, $d=B'(1)$ and $e=B''(1)$. Then
\begin{equation}
 D_{\rm cont}^2:=\operatorname{dist}(\tau,\mathcal V_{\rm cont})^2
 =2-2b^2+2bd-d^2-\frac{e^2}{4}.
 \label{eq:ns76-continuous-distance}
\end{equation}
This expression is nonnegative and at most every integer finite
squared error $E_N$. It vanishes exactly when RH holds.
The formula contains unknown off-critical zero information and
provides no unconditional upper bound tending to zero.
\end{proposition}
\begin{proof}
The normalized real dilations of $\sigma_1$ become the causal
translation semigroup applied to $F$. The Beurling--Lax theorem gives
that its closed generated subspace is the inner factor $B$ times
$\mathbb H^2$. This use of the full continuous semigroup is the reason
we do not replace it by the integer set without a separate argument.

For $\Re z>1/2$, let $k_z(s)=1/(s+\overline z-1)$ be the evaluation
kernel. Multiplication $M_B$ is an isometry and
$M_B^*k_z=\overline{B(z)}k_z$.
Differentiating at real $z=1$ is legitimate in the Hilbert norm,
as is seen directly from the integrable kernel derivatives. Since
$T=-\partial_z k_z|_1-\tfrac12\partial_z^2 k_z|_1$, it follows that
\[
 M_B^*T=-\frac{d+e/2}{s}+\frac{b+d}{s^2}-\frac{b}{s^3}.
\]
The complete boundary Gram for $(1/s,1/s^2,1/s^3)$ is
\[
 \begin{pmatrix}1&1&1\\1&2&3\\1&3&6\end{pmatrix};
\]
its entries follow by integrating the Laplace representatives
$e^{-t/2}t^{j-1}/(j-1)!$ on $t>0$.
Therefore $\|M_B^*T\|^2=2b^2-2bd+d^2+e^2/4$.
The orthogonal projection is $B M_B^*T$, with that same norm, and
$\|T\|^2=\|\tau\|^2=2$. This proves the distance identity and
nonnegativity without any sign assumption on $d$ or $e$.

The integer span is contained in the continuous span, giving
$E_N\geq D_{\rm cont}^2$. If RH holds, the product is empty and
$(b,d,e)=(1,0,0)$. Conversely, any off-critical zero $\rho$ gives a
zero of $B$ in this half-plane, whereas $T(\rho)\ne0$ since
$\rho\notin\{0,1\}$. Continuous evaluation in $\mathbb H^2$ then
prevents $T$ from belonging to the closed subspace $B\mathbb H^2$,
so its distance is positive.
\end{proof}

\begin{corollary}[The fixed component invisible to every block gain]
\label{cor:ns76-invisible-component}
Let $P_{\rm cont}$ and $P_N$ be the complete projections onto the
continuous and finite integer spaces. Then exactly
\[
 R_N=(I-P_{\rm cont})\tau+(P_{\rm cont}\tau-P_N\tau),
 \qquad
 E_N=D_{\rm cont}^2+\|P_{\rm cont}\tau-P_N\tau\|^2.
\]
Every new correlation $g_n$ and every selected-direction gain depend
only on the second component. The first is orthogonal to every real
dilation, hence also to each projected new integer atom.
\end{corollary}
\begin{proof}
The finite integer space is a subspace of $\mathcal V_{\rm cont}$,
so $P_NP_{\rm cont}=P_N$. The two displayed components are
orthogonal. Taking inner products with any old-projected new atom
annihilates the first, proving the final statement.
\end{proof}

\begin{proposition}[The integer and continuous closures are strictly different]
\label{prop:ns76-strict-closure}
Let $\mathcal V_{\rm int}$ be the closed span of the integer atoms.
Then $\mathcal V_{\rm int}\subsetneq\mathcal V_{\rm cont}$,
unconditionally. More explicitly,
\[
 \operatorname{dist}(\sigma_{3/2},\mathcal V_{\rm int})^2>.00111713.
\]
This is a bound for $\sigma_{3/2}$, not for the target $\tau$.
\end{proposition}
\begin{proof}
On $I=(1/2,1)$, the first integer atom is
$\sigma_1=1/x+\log x$ and every integer atom with $n\geq2$
is $1/(nx)$. Their restricted closed span lies in the closed
two-dimensional space $\operatorname{span}\{1/x,\log x\}$.
In contrast, $\sigma_u$ for $1<u<2$ has a nonzero derivative jump
at the interior point $x=1/u$; it cannot equal an element of that
analytic two-dimensional span, even almost everywhere. This already
proves strict inclusion.

For the quantitative bound, put $u=3/2$, $L=\log2$, $c=\log u$ and
$h=L-c$. The complete local Gram, loads and squared norm are
\[
 G_I=\begin{pmatrix}1&-L^2/2\\-L^2/2&1-L-L^2/2\end{pmatrix},\qquad
 b_I=\begin{pmatrix}
 u^{-1}-h^2/2\\
 (c+2-L^2/2)/u+c(L+1)/2-(L^2+2L+2)/2
 \end{pmatrix},
\]
\[
 q_I=u^{-2}-h^2/u+2/u-h^2/2-h-1.
\]
These follow by splitting only at $x=1/u$, with
$\sigma_u=1/(ux)+\log(ux)$ on the left and $1/(ux)$ on the right.
The Gram is positive definite. Its exact local projection error
$q_I-b_I^TG_I^{-1}b_I$ lies in $(.00111713,.00111714)$ by independent
256/384-bit enclosures; Maxima checks all complete elementary integrals.
Restriction of a norm to $I$ can only decrease it, so this local
minimum bounds the full-space distance below. No omitted tail is
set to zero in the asserted full-space inequality.
\end{proof}

\paragraph{What the calculation does and does not resolve.}
The $q=2$ continuous defect is now expressed through three values of the
classical unknown inner factor. It is not estimated away. In particular,
the integer limiting error may contain an additional approximation
component. The closures are strictly different, but that fact alone
does not decide whether their two distances from this particular target
are different. No equality of positive target distances is supplied.
The original sufficient cofinal contraction remains valid, because it
would force the entire $E_N$ to zero, including this fixed component.
Replacing $E_N$ by $E_N-D_{\rm cont}^2$ would remove exactly the part
that must be ruled out and would not prove RH.
